\section{Week 16: MINCORE KeyGen First Task}

The Week~16 first task originally exited as
\status{continue-kg: blocked-kg-ntt-mul}.  The continuation reported in the
next section supersedes that intermediate decision with
\status{stop-kg-ntt}, not \status{go-kg}.  A transparent
Hoare harness now composes the actual mode-2
\procname{_kp_m23_matrix} and \procname{_keypair_finalize_m23} procedures.
The faithful paper claims (KG-1)--(KG-4) and
\(A s = q j \pmod {2q}\) remain incomplete because the checked matrix output
has not yet been related to the security model's polynomial multiplication.

\subsection*{Baseline and actual boundary}

Before modification, the aggregate verifier fresh-compiled 74 authored
targets with \texttt{-no-eco}.  The preserved log is
\pathname{logs/verify-all-before-week16.log}, with SHA-256
\texttt{da7a7516166d54e93665d36e9a81348050ac6a83786111eb83ac6901743d6056}.

\theoryname{Mode2KeygenCoreEquation.ActualM23MatrixFinalizeSnapshot.run}
first invokes the generated parent procedure \procname{_kp_m23_matrix} with
\((\mathit{rows},\mathit{cols})=(2,3)\), records the returned
pre-finalization \texttt{bp} array, and then invokes
\procname{_keypair_finalize_m23} with count 512.  It returns the actual
snapshot, transformed-secret scratch, finalized high-part array, and adjusted
\texttt{s2}.  The harness contains no sampler, retry loop, packer, acceptance
guard, or public API.

\subsection*{Compiled snapshot consequences}

The Hoare precondition binds all six initial arrays and requires only the
existing matrix bound, the three input-polynomial representations, centered
active \texttt{s2} words, and canonical active \texttt{avec} words.  It does
not assume finalizer success, a KG predicate, or a desired equation.

\procname{actual_m23_matrix_finalize_snapshot} exports the actual matrix
output and NTT-scratch representations, exact finalizer output, and tail
frames.  The stronger
\procname{actual_m23_matrix_finalize_semantic_snapshot} additionally exports
the scalar and HAETAE finalizer semantics.  For every active coefficient, if
\[
 r=(\mathit{pre\_bp}+e+a)\bmod q,
 \qquad b_0=\mathsf{low}(r),
\]
then the compiled result includes
\[
 r=2b_1+b_0,\qquad -1\leq b_0\leq1,
 \qquad e'=e-b_0.
\]
This is the KG-2 coefficient decomposition and the adjusted component of
KG-4 at the actual finalizer boundary.

\theoryname{Mode2KeygenSnapshotAlgebra} proves the required integer division,
parity, and reduction lemmas internally.  The resulting actual predicate is
the deliberately snapshot-only congruence
\[
 2(a-2b_1)+2\,\mathit{pre\_bp}+2e'\equiv0\pmod {2q}.
\]
The name \procname{actual_snapshot_mod2q_zero} preserves that distinction;
the returned snapshot is not silently renamed to
\(A_{\mathrm{gen}}s_{\mathrm{gen}}\).

\subsection*{Exact blocker}

The checked matrix postcondition represents each returned row by
\[
 \mathsf{output\_row}
 =\mathsf{array256\_mont}
   (\mathsf{full\_invntt}(\mathsf{pointwise\_row\_ntt}(\cdots))).
\]
No compiled theorem identifies this object with the security model's row
product \(\sum_c A_{\mathrm{gen}}[r,c]s_{\mathrm{gen}}[c]\).  The legacy
\procname{Rq.ntt} cannot be substituted: the existing full-NTT specification
contains a theorem showing that the two transforms differ on one.

Consequently KG-1 and KG-3 are \status{blocked}; KG-4 is partial because only
the adjusted \(e-b_0\) component is constructed; and the paper
\(A s=qj\pmod {2q}\) conclusion is \status{blocked}.  The claim ledger keeps
\claimid{OBL-MINCORE-KEYGEN} partial.  Sampler distributions, retry
termination, packing, public APIs, Sign/Verify correctness, and security are
outside this result.

\subsection*{Historical continuation target}

The requested KeyGen continuation was one representation theorem: connect
\procname{KeygenM23ArithmeticSpec.output_row} to the security model's
\(A_{\mathrm{gen}}s_{\mathrm{gen}}\) polynomial product through the checked
full NTT.  The next section records why that bridge did not close and why the
active lane moves to Sign without assembling KG-1, KG-3, complete KG-4, or
the paper key equation.
